% Generated draft: requires mathematical and source review.
\documentclass[11pt]{article}
\usepackage[margin=25mm]{geometry}
\usepackage{fontspec}
\setmainfont{Latin Modern Roman}
\usepackage{amsmath,amssymb,mathtools,unicode-math}
\setmathfont{Latin Modern Math}
\usepackage{enumitem,needspace,xurl,hyperref}
\setlength{\parindent}{0pt}
\setlength{\parskip}{5pt}
\setlength{\emergencystretch}{3em}
\newcommand{\N}{\mathbb{N}}
\newcommand{\Z}{\mathbb{Z}}
\newcommand{\Q}{\mathbb{Q}}
\newcommand{\R}{\mathbb{R}}
\newcommand{\C}{\mathbb{C}}
\newcommand{\F}{\mathbb{F}}
\newcommand{\A}{\mathbb{A}}
\newcommand{\PP}{\mathbb{P}}
\newcommand{\E}{\mathbb{E}}
\newcommand{\eps}{\varepsilon}
\newcommand{\dd}{\,\mathrm{d}}
\begin{document}
\section*{UnsolvedMath ID: 800323}
\textbf{Source problem code:} catalogue entry
\subsection{Definitions and mathematical statement}
For each integer \(n\geq 1\), let \(p_n\) denote the \(n\)-th prime and let
\[
\pi(n)=\#\{p\leq n:p\text{ is prime}\}.
\]
For terminology, let \(K\) be a field of characteristic zero and let \(F(z)\in K[[z]]\). The series \(F\) is differentially transcendental over \(K(z)\) if, for every integer \(r\geq 0\) and every nonzero polynomial \(P\in K(z)[Y_0,\ldots,Y_r]\),
\[
P\bigl(F(z),F'(z),\ldots,F^{(r)}(z)\bigr)\neq 0.
\]
Here derivatives are formal derivatives and the expression is interpreted in \(K((z))\). The series is D-finite over \(K(z)\) if there exist an integer \(r\geq 0\) and polynomials \(a_0(z),\ldots,a_r(z)\in K[z]\), not all zero, such that
\[
\sum_{j=0}^{r}a_j(z)F^{(j)}(z)=0.
\]
These definitions explain the terminology; the supplied record does not specify the coefficient field for its intended assertion.

\textbf{Unresolved statement of Conjecture 7.} The supplied record attributes this conjecture to Pak (2018) and labels it open. Its extracted statement refers to two generating functions involving the preceding prime notation, but gives neither generating function. It ends with a malformed fragment followed by the words ``is D-transcendental.'' Thus the function or functions being conjectured to be differentially transcendental, their indexing conventions, and the base field cannot be recovered precisely from the supplied record. No definite conjectural formula can be stated without restoring the missing source text.

The supplied surrounding discussion attributes non-D-finiteness of both unidentified generating functions to Flajolet, Gerhold and Salvy, reference [FGS], using asymptotic arguments. It contrasts that reported result with the greater difficulty of proving differential transcendence and invokes the irrationality questions for \(e+\pi\) and \(\zeta(5)\) as analogies, with a reference to [KoZ]. These are supplied contextual assessments, not independently verified mathematical or literature-status claims.
\subsection{Short English statement}
Pak's Conjecture 7 concerns differential transcendence in a discussion of generating functions associated with primes. The supplied extraction omits the generating functions and the subject of the conjecture, so its exact mathematical statement requires source repair. The registry's open label and its reported non-D-finiteness results are supplied assessments requiring review.
\subsection{Sources}
\begin{itemize}
\item Local UnsolvedMath registry, permanent ID \(800323\), exact source problem code \texttt{catalogue entry}, title ``Complexity in Enumerative Combinatorics --- Conjecture 7.'' The supplied record attributes the problem to Pak (2018) and labels it open.
\item Pak, \emph{Complexity Problems in Enumerative Combinatorics} (2018), Conjecture 7. Supplied source URL: \url{https://arxiv.org/abs/1803.06636}. The source record reports access on 2026-10-16 and extraction by \texttt{arxiv\_latex}; these are supplied provenance details, not evidence of independent access or verification here.
\item Reference [FGS], attributed in the supplied text to Flajolet, Gerhold and Salvy, for an asymptotic proof of non-D-finiteness of both referenced generating functions. No full bibliographic entry or separate URL is supplied.
\item Reference [KoZ], cited in the supplied contextual discussion. No authors, title, bibliographic entry, or URL are supplied for this reference.
\end{itemize}
\end{document}
